% Generated draft: requires mathematical and source review.
\documentclass[11pt]{article}
\usepackage[margin=25mm]{geometry}
\usepackage{fontspec}
\setmainfont{Latin Modern Roman}
\usepackage{amsmath,amssymb,mathtools,unicode-math}
\setmathfont{Latin Modern Math}
\usepackage{enumitem,needspace,xurl,hyperref}
\setlength{\parindent}{0pt}
\setlength{\parskip}{5pt}
\setlength{\emergencystretch}{3em}
\newcommand{\N}{\mathbb{N}}
\newcommand{\Z}{\mathbb{Z}}
\newcommand{\Q}{\mathbb{Q}}
\newcommand{\R}{\mathbb{R}}
\newcommand{\C}{\mathbb{C}}
\newcommand{\F}{\mathbb{F}}
\newcommand{\A}{\mathbb{A}}
\newcommand{\PP}{\mathbb{P}}
\newcommand{\E}{\mathbb{E}}
\newcommand{\eps}{\varepsilon}
\newcommand{\dd}{\,\mathrm{d}}
\begin{document}
\section*{UnsolvedMath ID: 30006009}
\textbf{Source problem code:} catalogue entry
\subsection{Definitions and mathematical statement}
The supplied record attributes the following conjecture to Fargues and describes it as a slightly modified version of part of his conjecture. Its ambient hypotheses are incomplete: the local field \(E\), the class of groups \(G\), the coefficient prime \(\ell\), and the precise meaning of an \(L\)-parameter are not specified. The statement below retains these as unresolved source qualifications.

Write \(\widehat G\) for the dual group appearing in the source, and let \(\operatorname{Bun}_G\) denote the intended moduli stack of \(G\)-bundles in the geometrization of local Langlands. Write
\[
D_{\mathrm{lis}}(\operatorname{Bun}_G,\overline{\Q}_\ell)
\]
for the category used in the supplied clean statement. Its geometric base and categorical conventions must be supplied before this becomes a fully specified formulation. For a basic class \(b\in B(G)\), let \(G_b\) be the corresponding inner form, and write
\[
i_b:B G_b(E)\hookrightarrow\operatorname{Bun}_G
\]
for the open immersion asserted in the source. Here \(B G_b(E)\) denotes the classifying stack of \(G_b(E)\), whereas \(B(G)\) denotes the set indexing the bundle strata.

For a discrete \(L\)-parameter \(\phi\) of \(G\), let \(\Pi_\phi(G_b(E))\) denote the local Langlands packet predicted for this inner form. For \(\pi\) in this packet, write
\[
m_\phi(\pi)=\dim\rho_\pi,
\]
where \(\rho_\pi\) is the representation of the centralizer of \(\phi\) assigned to \(\pi\) by the intended local Langlands correspondence for extended pure inner forms. This is the multiplicity denoted \(\langle\pi,e\rangle\) in the source. The precise centralizer, enhancement convention, and auxiliary choices are not specified.

\textbf{Conjectural construction.} For every discrete \(L\)-parameter \(\phi\) of \(G\), construct an object
\[
F_\phi\in D_{\mathrm{lis}}(\operatorname{Bun}_G,\overline{\Q}_\ell)
\]
with the following properties.
\begin{enumerate}
\item \(F_\phi\) is a Hecke eigensheaf with eigenvalue \(\phi\). More explicitly, for every \(V\in\operatorname{Rep}(\widehat G)\), the intended eigenrelation has the schematic form
\[
T_V(F_\phi)\simeq \mathcal E_{\phi,V}\boxtimes F_\phi,
\]
where \(T_V\) is the Fargues--Scholze Hecke operator mentioned in the source and \(\mathcal E_{\phi,V}\) is the eigenvalue object obtained from \(\phi\) through \(V\). The isomorphisms must have the compatibilities required by the intended definition of a Hecke eigensheaf. This displayed relation does not fix the missing target category, parameter space, or normalization of \(T_V\).
\item \(F_\phi\) is perverse for the perverse \(t\)-structure referred to in the source as \([8,\text{ Proposition }1.2.1]\). The bibliographic identity of \([8]\) and the conventions for this \(t\)-structure are absent from the supplied record.
\item For every basic \(b\in B(G)\), there is an isomorphism
\[
i_b^*F_\phi\simeq
\bigoplus_{\pi\in\Pi_\phi(G_b(E))}
\pi^{\oplus m_\phi(\pi)}.
\]
The representations on the right are understood as the corresponding objects on \(B G_b(E)\). The source does not specify any normalization or cohomological shifts needed for this identification.
\end{enumerate}

The original extract also mentions intertwining \(T_V(-)\) with tensoring by an associated vector bundle on a parameter stack written as \(Z^1(W_E,\widehat G)_{\overline{\Q}_\ell}/\widehat G\). Its placement in the extracted text is corrupted, so the precise additional assertion and its relationship to the fixed-parameter eigenrelation above cannot be recovered reliably. Likewise, the domain and discreteness condition for \(\phi\) must be recovered from the source context; no choice between Weil, Weil--Deligne, or other parameter conventions is imposed here.
\subsection{Short English statement}
For each discrete local Langlands parameter, the conjecture asks for a perverse Hecke eigensheaf whose restriction to every basic stratum contains exactly the predicted packet representations, with multiplicities given by dimensions of their centralizer representations. The supplied extract leaves essential ambient hypotheses, normalizations, and part of the spectral compatibility assertion unresolved; this is a statement-only draft requiring human mathematical review.
\subsection{Sources}
\begin{itemize}
\item Local registry: supplied UnsolvedMath record, permanent ID \(30006009\), source problem code \texttt{catalogue entry}; project registry \texttt{data/unsolvedmath/problems.json}, with \texttt{catalog/problems.json} designated as the supplementary registry. These local files were not independently inspected for this draft.
\item Supplied primary report citation: \emph{Arithmetic Geometry (2025)}, Oberwolfach Reports, DOI \texttt{10.4171/owr/2024/33}, pp.\ 1855--1912. \url{https://doi.org/10.4171/owr/2024/33}. The record gives proposed year \(2024\), while its citation title gives \(2025\); this discrepancy is unresolved. The report is described as documenting mathematical workshops and their open-problem discussions.
\item Attribution: the original extract labels this a conjecture of Fargues, in a slightly modified version of part of his conjecture, and refers to an unidentified reference \([4]\). It also cites an unidentified \([8,\text{ Proposition }1.2.1]\) for perversity. The record supplies no proposer name in its proposer field.
\item Supplied literature source labeled ``Hamann--Imai--Yang'': \url{https://arxiv.org/abs/1608.07446}. The label and its connection to the stated claims have not been independently verified.
\item Supplied literature source labeled ``Fargues--Scholze'': \url{https://arxiv.org/abs/2102.13459}. Its connection to the stated claims has not been independently verified.
\item Supplied literature assessment dated \texttt{2026-08-22}: the record labels the problem ``partially solved'' and states that Hecke eigensheaf and packet properties are proved in important cases, including cuspidal parameters for general linear groups and special low-rank cases, while the general discrete-parameter conjecture remains open. This is supplied dated curation derived from the OWR material, not an independently verified literature assessment. The record expressly requests independent verification of time-sensitive status and cited claims.
\item The registry labels its statement \texttt{corrected\_verified} and says that only formula delimiters, notation, grammar, and typography were repaired. This supplied label does not establish mathematical correctness of the formulation or of this draft.
\end{itemize}
\end{document}
